
\section{平坦性的拓扑推论}

本节讨论平坦性的若干拓扑面向与推论。它们都可归结为平坦态射的 Going-Down 定理（\S24.5.3），而后者又源自忠实平坦性。由于代数几何中的维数是一个拓扑概念，我们将证明：在平坦\trFamily{}中，纤维的维数表现良好（\S24.5.5）。

\textbf{24.5.1. 忠实平坦。}\quad \trFaithfullyFlat{}性这个概念因种种缘由而颇为好用，这里只描述其中的几条。设 $M$ 是一个 $B$-\trModule。若对一切由 $B$-\trModule{}组成的\trComplexN
\[
\tag{24.5.1.1}\label{eq:faithful-flat-complex}
N' \longrightarrow N \longrightarrow N'',
\]
\eqref{eq:faithful-flat-complex} 正合当且仅当 \eqref{eq:faithful-flat-complex}$\otimes_B M$ 正合，我们就称 $M$ 是\textbf{\trFaithfullyFlat}的。若 $B$-\trAlgebra{} $A$ 作为 $B$-\trModule{}忠实平坦，我们就称 $A$ 是忠实平坦的。

\textbf{24.5.A.（练习）}\quad 证明：平坦 $B$-\trModule{} $M$ 是忠实平坦的，当且仅当对一切 $B$-\trModule{} $N$，$M\otimes_B N=0$ 蕴含 $N=0$。

\textbf{24.5.B.（练习）}\quad 设 $M$ 是平坦 $B$-\trModule。证明以下诸条等价。
（a）$M$ 忠实平坦；
（b）对一切\trPrimeIdeal{} $\mathfrak{p}\subset B$，$M\otimes_B\kappa(\mathfrak{p})$ 非零；
（c）对一切\trMaximalIdeal{} $\mathfrak{m}\subset B$，$M\otimes_B\kappa(\mathfrak{m})=M/\mathfrak{m}M$ 非零。

设 $\pi\colon X\to Y$ 是概形之间的一个\trMorphism。若 $\pi$ 既\trFlat{}又是\trSurjectiveMap，我们就称 $\pi$ 是\textbf{忠实平坦}的。（与平坦性不同，\trFaithfullyFlat{}性对\trQCSheaf{}而言并不特别有用，所以我们不定义基上的忠实平坦\trQCSheaf。）

\textbf{24.5.C.（练习（参看\mbox{练习~24.5.B}(b)））}\quad 设 $B\to A$ 是一个\trRingMorphism。证明：$A$ 在 $B$ 上忠实平坦，当且仅当 $\Spec A\to\Spec B$ 忠实平坦。

\textbf{24.5.2.}\ \emph{忠实平坦\trMorphism{}表现「合理」。}\quad 忠实平坦\trMorphism{}在 \S8.1 的意义下表现「\trReasonable」（在复合与基变换下保持，并且在靶上局部），因为平坦性与满射性都具有这些性质（分别见 \S24.1.8 与\mbox{练习~10.4.G}）。

\textbf{24.5.D.（练习）}\quad 设 $\pi\colon\Spec A\to\Spec B$ 是平坦的。
（a）证明：$\pi$ 忠实平坦，当且仅当每个\trClosedPoint{} $q\in\Spec B$ 都落在 $\pi$ 的\trImage{}中。（提示：\mbox{练习~24.5.B}(c)。）
（b）进而证明：局部环之间的\emph{任何}平坦（局部）\trMorphism（\mbox{定义~7.3.1}）都是忠实平坦的。（局部环之间的\trMorphism{}一律假定为局部的，即极大理想被拉回到极大理想。）

\textbf{24.5.3. 平坦态射的 Going-Down。}\quad \mbox{练习~24.5.D} 的一个推论，是下面这条有用的结果，它的陈述对\trFaithfullyFlat{}性只字未提。（这条陈述之所以令人想起\trIntegrallyClosed{}\trIntegralDomain{}的\trFinite\trExtension{}的 Going-Down 定理（\mbox{定理~12.2.12}），并非出于偶然。）

\textbf{24.5.E.（练习：平坦态射的 Going-Down 定理）}\quad
（a）设 $B\to A$ 是环的一个平坦\trMorphism，相应于映射 $\pi\colon\Spec A\to\Spec B$。设 $\mathfrak{q}\subset\mathfrak{q}'$ 是 $B$ 的\trPrimeIdeal，并设 $\mathfrak{p}'$ 是 $A$ 的\trPrimeIdeal，满足 $\pi([\mathfrak{p}'])=[\mathfrak{q}']$。证明：存在 $A$ 的\trPrimeIdeal{} $\mathfrak{p}\subset\mathfrak{p}'$，使得 $\pi([\mathfrak{p}])=[\mathfrak{q}]$。提示：证明 $B_{\mathfrak{q}'}\to A_{\mathfrak{p}'}$ 是平坦的局部环同态，从而忠实平坦（\mbox{练习~24.5.D}(b)）。\footnote{译注：原书此处作「flat by the Exercise 24.5.D(b)」，衍一冠词 the，当为「by Exercise 24.5.D(b)」之误，径改。——译者注}
（b）\mbox{（a）}给出了平坦性的一个几何推论。画一幅图来说明它。
（c）回想 \S8.2.4 中所描述的 Going-Up 定理。仿照\mbox{练习~8.2.F} 的方式陈述平坦态射的 Going-Down 定理，并证明之。
\textbf{24.5.F.（练习）}\quad 设 $\pi\colon X\to Y$ 是\trIntegral{}（比如\trFinite）的平坦\trMorphism，而 $Y$ 是\trPureDimensional{}的。证明 $X$ 是\trPureDimensional{}的。（这是对\mbox{练习~12.1.H}(a) 的推广。）提示：$\pi$ 同时满足 Going-Up（见\mbox{练习~8.2.F}）与 Going-Down（\mbox{练习~24.5.E}）。

\textbf{24.5.G.（重要练习：平坦\trMorphism{}是\trOpenMap{}（在合理的情形下））}\quad 设 $\pi\colon X\to Y$ 局部有限型且\trFlat，而 $Y$（从而 $X$）局部\trNoether。证明 $\pi$ 是\trOpenMap{}（即把\trOpenSet{}送到\trOpenSet）。提示：化归为证明：对每一个这样的 $\pi$，$\pi(X)$ 都是开的。再化归到 $X$ 与 $Y$ 都是仿射的情形。用\mbox{\trChevalley{}定理~8.4.2} 证明 $\pi(X)$ 是\trConstructible{}的。用平坦态射的 Going-Down 定理（\mbox{练习~24.5.E}）证明 $\pi(X)$ 在\trGenerization{}下稳定。最后用\mbox{练习~8.4.C} 得出结论。

\textbf{24.5.4.}\ \emph{\trhRemark。}\quad
（i）当然，开\trMorphism{}并不都是平坦的：$\Spec k[t]/(t)\to\Spec k[t]/(t^{2})$ 便是明证。
（ii）此外，在相当合理的情形下，平坦\trMorphism{}也\emph{不}必是开映射：$\Spec k(t)\to\Spec k[t]$ 便是明证（由\mbox{例~24.1.1}(ii) 它是\trFlat{}的）。

\textbf{24.5.H.}\ \hard\ \textbf{（练习）}\quad 把「局部有限型」与「局部诺特」这两个假设减弱成仅仅是「局部有限表现」。提示：正如\mbox{\trChevalley{}定理~8.4.2} 的类似推广（\mbox{练习~10.3.I}）中的做法，利用如下事实：任何这样的\trMorphism{}都是从某个\trNoether{}情势「局部地」拉回而来的。

\textbf{24.5.5. 对平坦\trMorphism{}而言纤维的维数表现良好。}

\textbf{24.5.6. \trhProposition。}\ ——\ \emph{设 $\pi\colon X\to Y$ 是局部\trNoether{}概形之间的一个平坦\trMorphism，$p\in X$ 与 $q\in Y$ 满足 $\pi(p)=q$。则}
\[
\operatorname{codim}_X p=\operatorname{codim}_Y q+\operatorname{codim}_{\pi^{-1}(q)}p
\]
（见\cref{fig:codim-fiber-inequality}）。

非正式的转述：纤维的维数就是 $X$ 与 $Y$ 的维数之差（至少局部地如此）。请与\mbox{练习~12.4.A} 对照：那里说的是，若没有平坦性假设，我们只会得到不等式（$\le$）。

\textbf{24.5.I.（练习）}\quad 按如下方式证明\mbox{命题~24.5.6}。正如刚才所说，\mbox{练习~12.4.A} 给出了一个方向的不等式，所以请证明另一个方向。给定 $Y$ 中含 $q$ 的一条\trIrreducible\trClosedSubset{}链，以及 $\pi^{-1}(q)\subset X$ 中含 $p$ 的一条\trIrreducible\trClosedSubset{}链，用平坦态射的 Going-Down 定理（\mbox{练习~24.5.E}）构造出 $X$ 中含 $p$ 的一条\trIrreducible\trClosedSubset{}链。

作为\mbox{命题~24.5.6} 的一个推论：若 $\pi\colon X\to Y$ 是不可约簇之间的平坦映射，那么 $\pi$ 的纤维都具有纯维数 $\dim X-\dim Y$。（\trhWarning{}：$\Spec k[t]/(t)\to\Spec k[t]/(t^{2})$ 并不展现纤维维数的跳变，它是开的，并且把\trAssociatedPoint{}送到\trAssociatedPoint——参看\mbox{练习~24.1.K}——但它并不平坦。如果你更想要一个既约的例子，\cref{fig:cusp-normalization} 所绘的尖点平面三次曲线的正规化 $\Spec k[t]\to\Spec k[x,y]/(y^{2}-x^{3})$ 也具有这些性质。）这便把我们引向下面这条有用的定义。

\textbf{24.5.7.}\ \emph{\trhDefinition。}\quad 若 $\pi\colon X\to Y$ 是局部有限型的平坦\trMorphism，而 $\pi$ 的一切纤维都具有纯维数 $n$，我们就说 $\pi$ 是\textbf{相对维数 $n$ 的平坦}（flat of relative dimension $n$）。（特别地，当人们说一个\trMorphism{}是相对维数 $n$ 的平坦\trMorphism{}时，「局部有限型」假设已蕴含其中。\mbox{注记~24.5.9} 为这一假设给出了动机。）

\textbf{24.5.J.（练习）}\quad 设 $\pi\colon X\to Y$ 是有限型 $k$-概形之间的一个平坦\trMorphism，而 $Y$ 是\trPureDimensional{}的（于是「\trCodimension{}就是\trDimension{}之差」；参看\mbox{定理~12.2.9}）。证明以下两条等价。
（i）概形 $X$ 具有纯维数 $\dim Y+n$。
（ii）\trMorphism{} $\pi$ 是相对维数 $n$ 的平坦\trMorphism。

\textbf{24.5.K.（练习）}\quad 证明：「相对维数 $n$ 的平坦\trMorphism」这一概念在任意基变换下保持。（为证纤维维数保持，提示：\mbox{练习~12.2.M}。）

\textbf{24.5.L.（练习）}\quad 设 $\pi\colon X\to Y$ 与 $\rho\colon Y\to Z$ 分别是相对维数 $m$ 与 $n$ 的平坦\trMorphism（从而是局部有限型的）。证明 $\rho\circ\pi$ 是相对维数 $m+n$ 的平坦\trMorphism。

\textbf{24.5.8.}\ \emph{\trhRemark。}\quad 作为练习，你现在可以为\mbox{练习~12.2.F}（维数对簇之积是可加的）给出第二个解法。

\textbf{24.5.9.}\ \emph{\trhRemark。}\quad 在「相对维数 $n$ 的平坦」的定义中要有「局部有限型」假设的缘由是：我们希望这是一个「\trReasonable」的\trMorphism{}类。特别地，我们希望「相对维数 $n$ 的平坦性」这一概念在基变换下保持。考虑\trCartesianDiagram
\[
\begin{tikzcd}[row sep=2.6em, column sep=3.4em]
\Spec k(x)\otimes_k k(y) \ar[r] \ar[d, "\pi'"] & \Spec k(y) \ar[d, "\pi"] \\
\Spec k(x) \ar[r] & \Spec k
\end{tikzcd}
\]
$\pi$ 与 $\pi'$ 都平凡地\trFlat，因为它们是到一些域的\trSpectrum{}的\trMorphism（\mbox{练习~24.1.A}(a)）。但 $\pi$ 的\trFiber{}维数为 $0$，而（如\mbox{注记~12.2.17} 所述）$\pi'$ 的\trFiber{}维数为 $1$。
\textbf{24.5.10. 到 $\Spec k$ 的\trMorphism{}是\trUniversallyOpen{}的。}\quad 我们现在来兑现 \S10.5.6 中的一个许诺。

\textbf{24.5.11. \trhTheorem。}\ ——\ \emph{每个\trMorphism{} $X\to\Spec k$ 都是\trUniversallyOpen{}的，即在\emph{任何}基变换之后仍为开。}

\hard\ \emph{证明。}\quad 我们希望证明：对任何 $k$-概形 $Y$，$X\times_k Y\to Y$ 都是\trTopologicalSpace{}之间的\trOpenMap，也就是说，它把\trOpenSet{}送到\trOpenSet。考虑 $Y$ 的一个\trAffine\trCover，便见可化归到 $Y$ 为仿射的情形，比如说 $\Spec S$。再考虑 $X$ 的一个\trAffine\trCover，便见又可化归到 $X$ 为仿射的情形，比如说 $\Spec R$。于是我们要证的是：$\pi\colon\Spec R\otimes_k S\to\Spec S$ 是\trOpenMap。

由于特异仿射开集构成仿射概形的\trZariskiTopology{}的一个\trBase，只须证明：对一切 $f\in R\otimes_k S$，$\pi(D(f))$ 都是开的。把 $f$ 写成有限和 $\sum_{i=1}^{n}r_i\otimes s_i$。设 $R'$ 是 $R$ 的由诸 $r_i$ 生成的\trSubring，于是 $R'\cong k[r_1,\dots,r_n]/I$，其中 $I$ 是某个（\trFinitelyGenerated{}的）\trIdeal。为免混淆，我们把 $f$ 在被视为 $R'\otimes_k S$ 中元素时记作 $f'$。考虑图
\[
\begin{tikzcd}[row sep=2.8em, column sep=4.2em]
\Spec R\otimes_k S \ar[rr, "\rho"] \ar[dr, "\pi"'] & & \Spec R'\otimes_k S \ar[dl, "\pi'"] \\
& \Spec S &
\end{tikzcd}
\]
由\mbox{练习~24.5.H}，$\pi'$ 是\trOpenMap（因为 $R'$ 在 $k$ 上\trFinitelyPresented），所以 $\pi'(D(f'))$ 是 $\Spec S$ 的一个\trOpenSubset。我们以证明 $\pi(D(f))=\pi'(D(f'))$ 来收尾。注意 $\rho^{-1}(D(f'))=D(f)$（因为 $\rho^{\flat}(f')=f$），于是我们有包含关系 $\pi(D(f))\subset\pi'(D(f'))$。为了证明反方向的包含，设 $p\in\pi'(D(f'))$。那么 $\pi'|_{D(f')}$ 在 $p$ 之上的\trFiber{}非空，也就是说，下述环不是\trZeroRing：
\[
(R'\otimes_k S)_{f'}\otimes_S\kappa(p)=(R'\otimes_k\kappa(p))_{f'}.
\]
于是 $R'\otimes_k\kappa(p)$ 不是零环，并且 $f'$ 在这个环中不是\trNilpotent{}的。但 $R'\otimes_k\kappa(p)$ 是 $R\otimes_k\kappa(p)$ 的一个\trSubring（$\cdot\otimes_k\kappa(p)$ 是\trExact{}函子；域上的任何东西都平坦），所以 $R\otimes_k\kappa(p)$ 也不是零环，并且 $f$ 在这个环中不\trNilpotent。于是 $(R\otimes_k\kappa(p))_f\neq 0$，从而 $\pi|_{D(f)}$ 在 $p$ 之上的\trFiber——即 $\Spec(R\otimes_k S)_f\otimes_S\kappa(p)$——也非空。\hfill$\square$

\textbf{24.5.12. 不加证明的几条更花哨的平坦性事实。}\quad 最后，我们以几条重要的事实收尾，但由于用不到它们，我们不予证明。
「一般平坦性」（generic flatness）其实叫做「泛平坦性」（general flatness）更贴切（参看\mbox{定义~3.6.11} 之后的讨论）。\footnote{译注：原文在这里玩了 generic（一般）与 general（泛）的文字游戏：既然「泛平坦性」更贴切，那么陈述自然也可以「泛」一些（in some generality）。}既然我们不给证明，那就不妨把它陈述得泛一些。\trNoether{}假设下一般平坦性的一个干净利落的证明，见 [Mu2, pp.~57--58]。

\textbf{24.5.13. \trhTheorem{}（一般平坦性，[Stacks, tag~052B]）。}\ ——\ \emph{设 $\pi\colon X\to Y$ 是概形之间的一个有限型\trMorphism，$Y$ 既约，而 $\mathscr{F}$ 是 $X$ 上的一个有限型\trQCSheaf。那么存在一个\trDense\trOpenSubscheme{} $\mathcal{U}\subset Y$，使得 $\mathscr{F}|_{\pi^{-1}(\mathcal{U})}$ 在 $\pi^{-1}(\mathcal{U})$ 上有限表现，并且在 $\mathcal{U}$ 上平坦。}

\textbf{24.5.14.}\ \emph{一般有限\trMorphism{}之\trDegree{}的解读。}\quad 我们将用它，借「数原像的个数」来解读簇的有理映射之次数，从而兑现 \S12.2.2 中的一个许诺。设 $\pi\colon X\dashrightarrow Y$ 是 $k$-簇之间的一个次数为 $d$ 的\trGenericallyFinite{}\trRationalMap。为简单起见，我们假设 $X$ 与 $Y$ 都\trIrreducible。（欢迎放宽这一点。）把 $X$ 换成使 $\pi$ 成为\trMorphism{}的那个开子集。由\mbox{命题~12.4.4}（「\trGenericallyFinite{}蕴涵泛有限」），存在 $Y$ 的一个\trDense\trOpenSubset{} $V$，使 $\pi$ 在其上\trFinite。由\mbox{定理~24.5.13}，存在 $V$ 的一个\trDense\trOpenSubset{} $V'$，使 $\pi$ 在其上\emph{\trFinite{}且\trFlat}。于是，由\mbox{注记~24.3.8}（「有限平坦\trMorphism{}具有局部常值的次数」），在这个开子集 $V'$ 上，$\pi$ 具有「局部常值的次数」，它必为 $d$。

\textbf{24.5.15.}\ \emph{\trFlat\trLocus{}是开的（平坦性是一个开条件）。}\quad 一般平坦性可以用来证明平坦性的最后一个拓扑面向：在合理的情形下，\trQCSheaf{}于一个基上平坦的轨迹是一个开子集。更精确地说：

\textbf{24.5.16. \trhTheorem{}（\trFlat\trLocus{}是开的）。}\ ——\ \emph{设 $\pi\colon X\to Y$ 是局部\trNoether{}概形之间的一个局部有限型\trMorphism，而 $\mathscr{F}$ 是 $X$ 上的一个\trCoherentSheaf。
（a）$X$ 中使 $\mathscr{F}$ 于 $Y$ 上平坦的那些点构成的\trLocus{}是 $X$ 的一个\trOpenSubset。
（b）若 $\pi$ 是闭的（比如\trProper），那么 $\mathscr{F}$ 在其上平坦的 $Y$ 中诸点构成的\trLocus（即那些使 $\mathscr{F}$ 在一切 $p\in\pi^{-1}(q)$ 处都平坦的 $q\in Y$）是 $Y$ 的一个\trOpenSubset。}

（b）可由（a）立即得到。（a）则化归为交换代数中一个非平凡的陈述；例如参看 [Mat2, Thm.~24.3] 或 [Gr-EGA, IV$_3$.11.1.1]。常有的情形是：\trNoether{}假设可以去掉，换成对\trMorphism{} $\pi$ 的局部有限表现假设；参看 [Gr-EGA, IV$_3$.11.3.1] 或 [Stacks, tag~00RC]。
